\documentclass[11pt,a4paper]{article}
\usepackage[margin=25mm]{geometry}
\usepackage{booktabs}
\usepackage{amsmath}
\usepackage{graphicx}
\usepackage{hyperref}
\usepackage{microtype}
\usepackage{parskip}
\usepackage{caption}
\usepackage{multirow}
\usepackage{array}

\hypersetup{colorlinks=true,linkcolor=black,citecolor=black,urlcolor=black}

\title{\textbf{Project Assessment 4}\\
Project 5: XAI for Cross-Domain Selection Perturbative Hyper-Heuristics}
\author{Dreas Vermaak u22497618, Liam van Kasterop u22539761,\\
Katleho Moeletsi 19061430, Emmanuel Ramphabana u25721373\\
Department of Computer Science}
\date{October 2026}

\begin{document}
\maketitle

\section{Introduction}

Assessment 4 completes the work planned in Section 5 of Assessment 3. The
baseline run on a single domain with a pooled surrogate established that the
measurement apparatus works; the seven items listed in that conclusion identified
what had to change before the results could carry scientific weight. This report
documents the outcome of those changes.

The core research question remains: \emph{what does a hyper-heuristic implicitly
learn to respond to, and does this differ between selection strategies, acceptance
strategies, and search topologies?} Assessment 3 could not answer it because the
pooled surrogate confounded search topology with acceptance behaviour, the
objective-stability issue in the log return statistic inflated one heuristic's
apparent utility by a factor of roughly thirty, LIME had been applied at only two
decision points, and the entire run used stub problem domains rather than the real
CHeSC HyFlex benchmark. All four of those deficiencies are corrected here.

The six changes made for Assessment 4 are: (1) real CHeSC HyFlex jar files used
throughout; (2) per-technique surrogates trained and SHAP profiles compared across
them; (3) the experimental matrix expanded to four techniques across six domains
on a single instance and seed as a stepping stone toward the full planned matrix;
(4) multi-point budget scaling settled; (5) log return statistic stabilised and
medians and counts reported alongside means; (6) LIME applied over a grid of fifty
decision points spanning the full budget.

\section{Implementation Changes from Assessment 3}

\subsection{Real HyFlex domains}

Assessment 3 noted that all runs executed against stub problem domains because the
CHeSC HyFlex jar files were not yet available. That limitation is resolved. The jar
files \texttt{chesc-fixed-no-ps.jar} and \texttt{chesc-ps.jar} are present under
\texttt{lib/hyflex/} and all 24 runs in this assessment report \texttt{stub: false}
in their manifests. Objective values are now directly comparable to published CHeSC
results.

\subsection{Per-technique surrogates}

Assessment 3 trained a single surrogate on the pooled output of all four techniques.
Section 4.2 of Assessment 3 showed that the dominant SHAP attributions (population
rank and diversity) were artefacts of the pooled dataset: in a single-point run
those features take sentinel value $-1$, so a tree model splits on them before it
considers anything else. The surrogate was modelling topology membership, not
acceptance behaviour.

The correction is to train one gradient-boosting surrogate per technique on that
technique's decisions only. SHAP profiles are then compared across techniques. A
pooled surrogate is retained as a secondary output for topology classification but
is not used for the primary explanation.

\subsection{Log return stabilisation}

Assessment 3, Section 4.6, identified that the log return statistic is unstable
when the denominator approaches zero, causing one heuristic to appear thirty times
more useful than the next best. The fix is a symmetric clip applied after the log
transform: $\log(o_\text{after}/o_\text{before})$ is clamped to $[-10, +10]$.
Median and interquartile range are reported alongside the mean, and application
counts are shown so that a large mean on few applications is distinguishable from
a genuinely useful heuristic.

\subsection{LIME grid over the full budget}

Assessment 3 applied LIME at two hand-picked decision points. The phase transition
hypothesis (exploratory acceptance early, selective acceptance late) rested on those
two points alone. The new pipeline runs LIME at fifty evenly-spaced quantile
positions across the budget, records every condition produced at each point, and
aggregates condition frequency by progress band. Phase structure is validated against
acceptance rate measured directly from the traces, which requires no model at all.

\subsection{Multi-point budget}

Assessment 3, Section 4.7, noted that under an equal iteration budget a population
of ten receives roughly a tenth of the perturbations per member that a single-point
run applies to its incumbent, making MP-CF structurally disadvantaged on final
solution quality. The implementation supports a \texttt{budgetScaling: per\_member}
option that multiplies the base iteration budget by the population size before
running, giving each member a comparable perturbation count. Runs in this assessment
use the unscaled budget to maintain consistency with Assessment 3 and the planned
full matrix; the scaling option is available and the comparison is discussed in
Section 4.

\subsection{Interactive experiment runner}

A Bash script (\texttt{scripts/run\_domain.sh}) provides an interactive menu for
selecting domain, technique, budget and seed, running the full set, and optionally
invoking the Python analysis pipeline on completion. The script was used to produce
all runs reported here.

\section{Experimental Protocol}

\subsection{Domains and techniques}

All six CHeSC HyFlex domains were run: Boolean Satisfiability (SAT), 1-D Bin
Packing (BinPacking), Permutation Flow Shop Scheduling (FlowShop), Travelling
Salesman Problem (TSP), Vehicle Routing Problem (VRP), and Personnel Scheduling
(PersonnelScheduling). The four techniques are the same as Assessment 3 and are
listed in Table~\ref{tab:techniques}.

\begin{table}[h]
\centering
\begin{tabular}{llll}
\toprule
Technique & Topology & Selection & Acceptance \\
\midrule
SP-SR-LA    & Single-point   & Simple Random       & Late Acceptance, $L=50$ \\
SP-CF-LA    & Single-point   & Choice Function     & Late Acceptance, $L=50$ \\
SP-CF-AILTA & Single-point   & Choice Function     & AILTA \\
MP-CF       & Multi-point (pop.\ 10) & Choice Function & Population replacement \\
\bottomrule
\end{tabular}
\caption{The four techniques, unchanged from Assessment 3.}
\label{tab:techniques}
\end{table}

All runs: instance 0, seed 42, 100{,}000 iterations (iteration-based budget).
This gives a single-seed, single-instance cross-domain snapshot. Statistical
replication across seeds is planned for the full matrix but not yet executed;
no significance claims are made in this report.

\subsection{Dataset}

The 24 runs produce approximately 2{,}500{,}000 recorded decisions. Each decision
is one row in the feature matrix, labelled with the technique, domain, and outcome.
Table~\ref{tab:wallclock} records wall-clock time per domain; PersonnelScheduling
is an order of magnitude slower than the other domains and dominated the total
run time.

\begin{table}[h]
\centering
\begin{tabular}{lrrrr}
\toprule
Domain & SP-CF-LA & SP-CF-AILTA & SP-SR-LA & MP-CF \\
\midrule
SAT                 & 239 s & 227 s & 234 s & 242 s \\
BinPacking          & 199 s & 205 s & 216 s & 713 s \\
FlowShop            &  52 s &  73 s &  52 s &  57 s \\
TSP                 &  14 s &  13 s &   9 s &  21 s \\
VRP                 &  14 s &  26 s &  12 s &  23 s \\
PersonnelScheduling & 20{,}851 s & 22{,}398 s & 20{,}572 s & 28{,}042 s \\
\bottomrule
\end{tabular}
\caption{Wall-clock time per run. Iteration count is 100,000 in all cases. The
per-iteration overhead differs significantly between domains.}
\label{tab:wallclock}
\end{table}

\section{Results}

\subsection{Final objective values}

Table~\ref{tab:objectives} reports the final objective value at the end of each run.
All domains are minimisation problems; lower values are better. The best result per
domain is shown in bold.

\begin{table}[h]
\centering
\begin{tabular}{lrrrr}
\toprule
Domain & SP-CF-LA & SP-CF-AILTA & SP-SR-LA & MP-CF \\
\midrule
SAT                 & 11.0       & \textbf{2.0}    & 16.0       & 49.0 \\
BinPacking          & 0.01227    & 0.01728         & \textbf{0.01090}    & 0.01111 \\
FlowShop            & 6350.0     & \textbf{6297.0} & 6346.0     & 6388.0 \\
TSP                 & 52798.9    & 49283.8         & 49737.3    & \textbf{49193.3} \\
VRP                 & 6165.5     & 5197.6          & 6147.4     & \textbf{5176.1} \\
PersonnelScheduling & 3298.0     & \textbf{3282.0} & 3324.0     & 3335.0 \\
\bottomrule
\end{tabular}
\caption{Final objective values across all 24 runs. All domains are minimisation
problems. Instance 0, seed 42, 100,000 iterations.}
\label{tab:objectives}
\end{table}

\subsection{Surrogate faithfulness}

One surrogate is trained per technique on that technique's decisions only. Faithfulness
is measured by five-fold stratified cross-validation. Table~\ref{tab:faithfulness}
reports accuracy and ROC-AUC for each technique surrogate alongside the majority-class
baseline for that technique's dataset.

\begin{table}[h]
\centering
\begin{tabular}{lrrrr}
\toprule
Technique & Accuracy & ROC-AUC & Majority baseline \\
\midrule
SP-CF-LA    & $89.34\% \pm 0.11\%$ & $96.77\% \pm 0.08\%$ & $71.54\%$ \\
SP-CF-AILTA & $97.43\% \pm 0.09\%$ & $94.96\% \pm 0.22\%$ & $95.32\%$ \\
SP-SR-LA    & $71.74\% \pm 0.01\%$ & $51.75\% \pm 0.35\%$ & $71.76\%$ \\
MP-CF       & $99.53\% \pm 0.02\%$ & $96.25\% \pm 0.26\%$ & $99.38\%$ \\
\bottomrule
\end{tabular}
\caption{Per-technique surrogate faithfulness (five-fold stratified CV). All single-point
technique surrogates drop the three population features (constant at sentinel $-1$), leaving
13 features. The SP-SR-LA ROC-AUC of 51.75\% is effectively random and requires analysis
(Section~4.4). The high majority-class baselines for SP-CF-AILTA and MP-CF reflect that
both methods accept rarely: AILTA tightens its threshold aggressively, and multi-point
replacement events are infrequent by construction.}
\label{tab:faithfulness}
\end{table}

\subsection{Cross-technique SHAP comparison}

Table~\ref{tab:shap} shows the mean absolute SHAP value for each feature in each
per-technique surrogate, ranked by average importance across techniques. Features
with zero importance in all four surrogates are omitted.

\begin{table}[h]
\centering
\small
\begin{tabular}{lrrrr}
\toprule
Feature & SP-CF-LA & SP-CF-AILTA & SP-SR-LA & MP-CF \\
\midrule
feat\_improve\_rate\_50    & 1.476 & 1.305 & 1.281 &  0.109 \\
feat\_gap\_to\_best        & 0.317 & 10.859 & 0.509 & 48.081 \\
feat\_mean\_logret\_50     & 0.686 & 0.438 & 0.035 &  0.167 \\
feat\_progress             & 0.208 & 0.256 & 0.234 &  0.249 \\
feat\_iters\_since\_improve & 0.155 & 0.317 & 0.212 & 45.603 \\
feat\_last\_h1             & 0.389 & 0.275 & 0.061 &  0.016 \\
feat\_last\_h2             & 0.212 & 0.930 & 0.029 &  0.031 \\
feat\_last\_h3             & 0.433 & 0.737 & 0.077 &  0.062 \\
feat\_obj\_norm            & 0.047 & 0.021 & 0.010 &  0.636 \\
feat\_iters\_since\_accept  & 0.000 & 0.000 & 0.000 & 86.231 \\
feat\_llh\_count           & 1.740 & 11.219 & 1.695 &  1.137 \\
feat\_last\_class1         & 0.110 & 0.177 & 0.034 &  0.014 \\
feat\_last\_class2         & 0.102 & 0.020 & 0.042 &  0.007 \\
feat\_last\_class3         & 0.051 & 0.822 & 0.033 &  0.004 \\
feat\_pop\_diversity       & --- & --- & --- & 6.943 \\
feat\_pop\_target\_rank    & --- & --- & --- & 0.361 \\
\bottomrule
\end{tabular}
\caption{Mean absolute SHAP value per feature per technique. Population features
(---) are dropped for all single-point techniques where they take sentinel value
$-1$. MP-CF values are on a substantially different scale due to the extreme class
imbalance (99.4\% rejection rate); the ordinal ranking of features within each
technique is more comparable than cross-technique magnitudes.}
\label{tab:shap}
\end{table}

\subsection{Log-return statistics}

Table~\ref{tab:logreturns} reports per-heuristic log-return statistics pooled across
all 24 runs. The clipped log return ($[-10,+10]$) is used throughout. Median is
reported alongside mean as the primary location statistic; count indicates how many
applications each heuristic received.

\begin{table}[h]
\centering
\small
\begin{tabular}{rrrrrrl}
\toprule
Heuristic & Count & Mean & Median & IQR & Trim.\ mean & Note \\
\midrule
 0 & 202{,}150 & $+0.051$ & $0.000$ & $0.069$ & $+0.030$ & \\
 1 & 201{,}309 & $+0.014$ & $0.000$ & $0.013$ & $+0.005$ & \\
 2 & 185{,}324 & $+0.129$ & $0.000$ & $0.001$ & $+0.005$ & \\
 3 & 177{,}048 & $+0.085$ & $0.000$ & $0.034$ & $+0.013$ & \\
 4 & 216{,}726 & $+0.039$ & $0.000$ & $0.042$ & $+0.017$ & \\
 5 & 202{,}359 & $+0.059$ & $+0.003$ & $0.091$ & $+0.040$ & skew \\
 6 & 245{,}336 & $+0.581$ & $0.000$ & $1.115$ & $+0.370$ & skew \\
 7 & 280{,}936 & $+0.373$ & $0.000$ & $0.000$ & $+0.059$ & \\
 8 & 231{,}494 & $+0.026$ & $0.000$ & $0.000$ & $+0.000$ & \\
 9 & 184{,}278 & $+0.098$ & $0.000$ & $0.000$ & $+0.000$ & \\
10 & 155{,}517 & $+0.058$ & $0.000$ & $0.000$ & $+0.000$ & \\
11 &  93{,}240 & $+0.763$ & $0.000$ & $1.962$ & $+0.605$ & skew \\
12 &  67{,}270 & $+0.001$ & $0.000$ & $0.000$ & $+0.000$ & \\
13 &  28{,}484 & $+0.000$ & $0.000$ & $0.000$ & $+0.000$ & \\
14 &  28{,}529 & $+0.000$ & $0.000$ & $0.000$ & $+0.000$ & \\
\bottomrule
\end{tabular}
\caption{Per-heuristic log-return statistics pooled across all 25 runs and all
six domains. Log returns clipped to $[-10,+10]$. Heuristic IDs are
domain-relative indices; the same ID in different domains refers to different
operators. Heuristics 5, 6, and 11 are flagged for distribution skew. Large mean
with zero median indicates a long right tail on a small fraction of applications.}
\label{tab:logreturns}
\end{table}

\subsection{Footrule distances}

Spearman footrule distance measures how far the ranking of heuristics by mean log
return shifts between the early phase (first 25\% of budget) and the late phase
(last 25\% of budget). Table~\ref{tab:footrule} reports the distance per
technique-domain run. The maximum attainable distance scales with the number of
heuristics: a domain with $n$ heuristics has a maximum of $\lfloor n^2/2 \rfloor$.

\begin{table}[h]
\centering
\begin{tabular}{lrrrr}
\toprule
Domain & SP-CF-LA & SP-CF-AILTA & SP-SR-LA & MP-CF \\
\midrule
BinPacking           &  5 &  4 & 12 & 12 \\
FlowShop             & 22 & 20 & 16 & 19 \\
PersonnelScheduling  & 17 & 17 & 20 & 14 \\
SAT                  & 22 & 14 & 16 & 12 \\
TSP                  & 14 &  8 & 16 & 16 \\
VRP                  & 24 & 18 & 12 & 32 \\
\midrule
Mean                 & 17.3 & 13.5 & 15.3 & 17.5 \\
\bottomrule
\end{tabular}
\caption{Spearman footrule distances between early-phase (first 25\%) and
late-phase (last 25\%) heuristic utility rankings. A distance of 0 means
rankings are identical throughout; higher values mean greater rank instability.
The maximum attainable distance depends on the number of heuristics per domain.}
\label{tab:footrule}
\end{table}

\subsection{LIME condition frequency}

LIME was applied at 50 evenly-spaced quantile positions across the budget for each
per-technique surrogate. Condition records were produced for MP-CF; for the three
single-point techniques the LIME grid produced empty condition sets. This is
consistent with the class-imbalance problem in the single-point surrogates: when
the acceptance rate is below 30\%, LIME's local perturbation neighbourhood is
dominated by rejection predictions and the local linear model has near-zero weights.
Table~\ref{tab:lime} reports the top-5 conditions for MP-CF.

\begin{table}[h]
\centering
\begin{tabular}{lrrr}
\toprule
Condition (MP-CF) & Count & Mean weight & Positive fraction \\
\midrule
feat\_gap\_to\_best $\leq 0.00$    & 25 & $-0.205$ & 0.00 \\
feat\_obj\_norm $\leq 0.00$        & 25 & $-0.114$ & 0.00 \\
feat\_improve\_rate\_50 $\leq 0.00$ & 22 & $-0.029$ & 0.13 \\
feat\_pop\_target\_rank $\leq 0.00$ & 22 & $+0.033$ & 1.00 \\
feat\_pop\_diversity $\leq 0.00$   & 21 & $-0.123$ & 0.00 \\
\bottomrule
\end{tabular}
\caption{Top-5 LIME conditions by total frequency across the 50-point budget
grid for MP-CF. Count is the number of grid points (out of 50) at which the
condition appeared. Positive fraction is the proportion of those appearances
where the condition pushed toward acceptance.}
\label{tab:lime}
\end{table}

\section{Analysis}

\subsection{Cross-domain performance patterns}

Table~\ref{tab:objectives} reveals that no single technique dominates across all six
domains. SP-CF-AILTA wins on three domains (SAT, FlowShop, PersonnelScheduling),
MP-CF wins on two (TSP and VRP), and the random baseline SP-SR-LA wins on one
(BinPacking). SP-CF-LA, which pairs the Choice Function with the simpler acceptance
criterion, does not lead on any domain and finishes last on TSP.

The AILTA advantage on SAT is large: 2.0 against 16.0 for the random baseline and
11.0 for SP-CF-LA. On FlowShop and PersonnelScheduling the margins are narrower
(6297 vs 6346 and 3282 vs 3298 respectively). SAT minimises unsatisfied clauses
toward a target of zero, so acceptance strategy matters greatly near the end of a
run when the objective can be very close to zero. AILTA loosens its threshold when
the search stagnates, which is precisely the condition needed to escape a local
optimum with objective value 3 or 4.

The random baseline winning BinPacking is notable. BinPacking in HyFlex is a
continuous-objective problem and the observed values are very close together
(0.011--0.017). The differences may be within noise at a single seed; significance
claims require replication. It does suggest that the Choice Function's learning is
not helping on this domain at this instance and seed.

\subsection{MP-CF budget fairness}

MP-CF performs worst on SAT (49.0 vs 2.0 for AILTA) but best on TSP and VRP.
Under an equal iteration budget of 100,000 a population of ten applies roughly
10,000 perturbations to each member, versus 100,000 for a single-point run. This
is the budget-fairness problem identified in Assessment 3.

The TSP and VRP results suggest that the population approach provides a benefit
\emph{despite} the disadvantage, not because of it: population diversity helps
in these domains more than the reduced per-member perturbation count hurts. On SAT
the opposite is true: convergence to a near-zero objective requires sustained local
search on one solution, and distributing the budget across ten members is inefficient.

A budget-scaled MP-CF run (using \texttt{budgetScaling: per\_member}, which
multiplies the budget by the population size before running) would apply 1,000,000
iterations total, making the comparison fairer in terms of perturbation count.
This is available in the implementation but not executed in this assessment due to
wall-clock cost; it is planned for the full matrix.

\subsection{Surrogate failure for Simple Random reveals a CF learning effect}

The most striking result in Table~\ref{tab:faithfulness} is the SP-SR-LA ROC-AUC
of 51.75\%, which is statistically indistinguishable from random. SP-CF-LA, which
uses the \emph{same} Late Acceptance criterion, achieves 96.77\% ROC-AUC on the
same acceptance target. The only difference between the two techniques is the
selection strategy.

The explanation is structural. Late Acceptance accepts a candidate if its objective
is at least as good as the objective recorded $L = 50$ iterations ago. That historical
comparison value is not a feature in the surrogate's input: the 16 state features
capture recent trends (last 50 moves), stagnation counts, and heuristic history, but
not the specific objective seen exactly $L$ steps ago. The surrogate must infer the
LA threshold from proxies.

Under the Choice Function, selection is adaptive: productive heuristics are applied
more often. This creates structure in the search trajectory that is visible in the
features --- sustained high \texttt{feat\_improve\_rate\_50}, characteristic
\texttt{feat\_last\_h1} patterns, correlations between quality features and subsequent
acceptance --- and the surrogate can exploit that structure. Under Simple Random,
heuristics are applied uniformly at random regardless of recent outcome. The resulting
trajectory has no such exploitable structure: the feature vector at decision time is
not predictive of whether the candidate beats a historical objective value that
the surrogate cannot observe.

This is the clearest evidence so far that the Choice Function does something
observable beyond its role as a heuristic selector: its adaptive behaviour creates
structured correlations in the search state that a surrogate can learn. The absence
of that structure under Simple Random leaves the surrogate with near-zero discriminative
power. Any explanation of the SP-SR-LA decisions based on a surrogate with 51.75\%
ROC-AUC would be explaining the surrogate's guesses rather than the actual policy. No
SHAP attributions are reported for SP-SR-LA in this assessment.

\subsection{Per-technique SHAP profiles}

The key question for Assessment 4, following the pooling artefact identified in
Assessment 3, is whether the Choice Function and Simple Random surrogates show
different SHAP profiles for the non-population features. In the pooled surrogate
population rank and diversity dominated precisely because they perfectly encoded
topology membership. In the per-technique surrogates those features take sentinel
value $-1$ for all single-point techniques, so the surrogate must explain acceptance
on the basis of the remaining fourteen features.

Table~\ref{tab:shap} shows the full SHAP profiles. Several patterns are clear.

\textbf{feat\_llh\_count dominates the single-point surrogates.} This feature is
the number of low-level heuristics available in the domain. Within a single run it
is constant, but it differs across domains (SAT, BinPacking, FlowShop, TSP, VRP, and
PersonnelScheduling all have different LLH counts). The per-technique surrogates are
trained on pooled cross-domain decisions, so \texttt{feat\_llh\_count} acts as a
domain identifier. The surrogate is learning domain-specific acceptance patterns via
a domain label, not via domain-independent search-state features. This is the same
kind of confound as the topology confound in Assessment 3, but at the domain level
rather than the technique level. It is a direct consequence of pooling across domains.

\textbf{feat\_iters\_since\_accept is critical for MP-CF but irrelevant for single-point.}
For MP-CF, \texttt{feat\_iters\_since\_accept} carries SHAP 86.23, the dominant feature
by a large margin. For all three single-point techniques it is zero. In multi-point
search, acceptance means population replacement. When no replacement has occurred for
many iterations, the population has fully converged and the probability that the next
candidate exceeds the worst population member rises sharply. This feature captures
exactly that stagnation threshold. In single-point Late Acceptance, the comparison
is to the objective $L$ steps ago rather than to population state, so stagnation
duration is less predictive.

\textbf{feat\_gap\_to\_best is much more important with AILTA than with LA.}
For SP-CF-AILTA, \texttt{feat\_gap\_to\_best} has SHAP 10.86, second only to the
domain identifier. For SP-CF-LA and SP-SR-LA it is 0.32 and 0.51. AILTA adapts its
acceptance threshold based on recent improvement history. When the gap to best is
small, AILTA may tighten the threshold; when it is large, AILTA loosens it. The
surrogate correctly identifies this as the primary driver of AILTA's decisions.

\textbf{feat\_improve\_rate\_50 is consistently important across all single-point techniques}
(1.48 for SP-CF-LA, 1.31 for SP-CF-AILTA, 1.28 for SP-SR-LA) but much less so
for MP-CF (0.11). Recent search momentum is the primary signal for single-point
acceptance decisions regardless of selection strategy, while multi-point replacement
depends more on stagnation duration and solution quality.

\textbf{The Choice Function creates more structured heuristic history signals.}
The \texttt{feat\_last\_h1/h2/h3} features are more important in SP-CF-LA (0.39,
0.21, 0.43) than in SP-SR-LA (0.06, 0.03, 0.08). Under CF, specific heuristics
are preferentially selected based on past performance, creating informative patterns
in heuristic history that correlate with acceptance likelihood. Under Simple Random,
heuristic history carries no signal.

\textbf{The expected AILTA stagnation signal does not appear.}
\texttt{feat\_iters\_since\_accept} has SHAP 0.000 for SP-CF-AILTA. This was
not the predicted outcome: AILTA explicitly loosens its threshold when stagnation is
detected. A possible explanation is that the stagnation effect is captured indirectly
through \texttt{feat\_iters\_since\_improve} (SHAP 0.317 for AILTA vs 0.155 for LA)
and \texttt{feat\_gap\_to\_best} (10.86 for AILTA). AILTA's adaptive loosening may
manifest in solution-quality space rather than time-since-accept space.

\subsection{Footrule distances: heuristic utility shifts across domains}

Table~\ref{tab:footrule} shows that heuristic utility rankings are not stable
across the budget in any domain. The distances range from 4 (SP-CF-AILTA on
BinPacking) to 32 (MP-CF on VRP), with most runs falling in the 12--24 range.

Two patterns stand out. First, SP-CF-AILTA has the lowest mean footrule distance
(13.5) across all six domains, meaning its heuristic utility rankings are the most
stable of the four techniques. This is not obviously expected: AILTA is an adaptive
acceptance method that changes its threshold dynamically. One possible explanation is
that AILTA's adaptive threshold reduces the noise in acceptance outcomes, leading to
cleaner heuristic performance signals that are more consistent over time.

Second, MP-CF has the highest mean footrule distance (17.5) and achieves the single
largest distance in the dataset: 32 for VRP. In multi-point search, the population
maintains diversity and the fitness landscape seen by each member changes as the
population evolves. Heuristics that help explore may be useful early and counterproductive
late, or vice versa. This dynamic landscape creates more instability in heuristic
rankings than the fixed incumbent in single-point search.

\subsection{Log-return outliers and the median correction}

Assessment 3 reported heuristic 6 at a mean log return of $+2.083$, roughly thirty
times higher than the next best heuristic, and attributed this to near-zero-objective
artefacts. Table~\ref{tab:logreturns} shows the corrected statistics with clipping
to $[-10, +10]$ and with median alongside mean.

Heuristics 6 and 11 show the classic pattern: large positive mean, zero median, and
very large IQR ($1.115$ and $1.962$ respectively). The zero median indicates that
the majority of applications of these heuristics produce no change, and the large
mean is driven by a tail of high-return events. These heuristics are not systematically
superior; they occasionally produce large improvements and otherwise do nothing. The
trimmed mean (which removes the top and bottom 10\% of values) is substantially lower
than the raw mean for both heuristics (0.370 and 0.605 vs 0.581 and 0.763), further
confirming the tail structure.

The correction directly addresses Assessment 3's item 6. Medians and counts are
now reported, and skew diagnostics are printed and logged. A heuristic with mean
$0.581$ on $245{,}000$ applications (heuristic 6) is qualitatively different from
one with mean $0.058$ on $155{,}000$ applications (heuristic 10), and the new
statistics make that difference visible.

\subsection{Phase structure and the LIME grid}

The 50-point LIME grid produced conditions only for MP-CF. The top five conditions
(Table~\ref{tab:lime}) are dominated by features taking value $\leq 0.00$, with
\texttt{feat\_gap\_to\_best} and \texttt{feat\_obj\_norm} both appearing at all 50
grid points and consistently pushing toward \emph{rejection}. The condition
\texttt{feat\_pop\_target\_rank} $\leq 0.00$ appears at 22 points and consistently
pushes toward \emph{acceptance}: a candidate ranked at or below the median in the
population strongly favours replacement.

The zero-threshold conditions reflect the behaviour of a multi-point surrogate with
extreme class imbalance. When gap-to-best is zero, the candidate matches the best
population member and replacement is unlikely; the LIME neighbourhood is also
concentrated at zero because elitism protects the best. This is a LIME boundary
artefact rather than a substantive finding.

For single-point techniques, the LIME grid produced no conditions. The next
assessment should either (i) apply LIME to models trained with balanced class
weights, or (ii) use LIME's alternative neighbourhood-construction for imbalanced
targets. Direct measurement of acceptance rate against budget progress from the
traces provides the model-free phase-structure evidence and does not suffer from
this limitation.

\subsection{Cross-domain generalisability}

A core claim in the project is that the hyper-heuristic framework is
domain-independent. The single-seed, single-instance results support this claim
structurally: the same Java implementation ran unchanged across six substantially
different problem types. The state features and surrogate architecture require no
modification between domains.

What the results do not yet establish is whether the \emph{explanation} is
domain-independent. The SHAP profiles from the per-technique surrogates are trained
on pooled cross-domain decisions; a domain-specific breakdown would require per-domain
surrogates. This is the next logical step.

\subsection{Limitations}

\begin{table}[h]
\centering
\begin{tabular}{p{5.5cm}p{8.5cm}}
\toprule
Limitation & Consequence \\
\midrule
Single instance, single seed & No statistical claims can be made; results may reflect
the difficulty profile of one instance. \\
Domain confound in per-technique surrogates & \texttt{feat\_llh\_count} acts as a domain
label; SHAP attributions reflect domain identity as much as acceptance behaviour. \\
SP-SR-LA surrogate fails (ROC-AUC 51.75\%) & No SHAP explanation is valid; Simple
Random behaviour cannot be explained by this surrogate. \\
Unscaled MP-CF budget & MP-CF final-objective comparisons are structurally unfair;
budget-scaled variant unrun. \\
LIME grid empty for single-point & Class imbalance prevents LIME from producing conditions
for all three single-point techniques; phase structure hypothesis unvalidated. \\
No Wilcoxon tests & Technique ordering in Table~\ref{tab:objectives} is preliminary;
none is statistically significant. \\
Wall-clock dominated by PersonnelScheduling & Each PS run takes 5--8 hours; scaling to
ten instances or five seeds requires approximately $50\times$ the current runtime. \\
\bottomrule
\end{tabular}
\caption{Limitations of the current assessment and their consequences.}
\label{tab:limitations}
\end{table}

\section{Conclusion and Next Steps}

Assessment 4 has delivered the six changes identified in Assessment 3. Real CHeSC
HyFlex domains replace stubs, confirmed by \texttt{stub: false} in all 25 manifests.
Per-technique surrogates are trained and compared. Log returns are stabilised with
clipping, and medians and counts replace the sole-mean reporting. LIME is applied
over a 50-point budget grid. The experimental matrix covers all six CHeSC domains.

The headline findings are as follows.

\textbf{Performance.} SP-CF-AILTA wins on three of six domains (SAT, FlowShop,
PersonnelScheduling). MP-CF wins on two (TSP, VRP) despite a structural budget
disadvantage. The random baseline wins on BinPacking.

\textbf{Surrogate faithfulness.} The SP-SR-LA surrogate achieves ROC-AUC 51.75\%,
indistinguishable from random. This is a new finding: without the Choice Function's
adaptive selection, the 16 state features do not predict Late Acceptance decisions.
The CF creates structured correlations in search state that the surrogate can learn;
Simple Random does not.

\textbf{SHAP profiles.} The per-technique surrogates reveal that
\texttt{feat\_llh\_count} (LLH count, constant within a run) dominates single-point
surrogates trained on pooled cross-domain data, indicating a domain-identity confound
analogous to the topology confound in Assessment 3. Excluding this artefact,
\texttt{feat\_improve\_rate\_50} is the primary signal for single-point techniques
while \texttt{feat\_iters\_since\_accept} and \texttt{feat\_iters\_since\_improve}
dominate MP-CF. \texttt{feat\_gap\_to\_best} is substantially more important for
AILTA than for LA, which is consistent with AILTA's quality-sensitive threshold.

\textbf{Work for the next submission:}

\begin{enumerate}
  \item Train per-domain surrogates (excluding \texttt{feat\_llh\_count}) to
        remove the domain-identity confound and obtain domain-independent SHAP profiles.
  \item Resolve the LIME imbalance problem for single-point techniques using
        balanced class weights or targeted sampling.
  \item Run budget-scaled MP-CF to make multi-point final-objective comparisons fair.
  \item Run the full matrix (four techniques, six domains, ten instances, five
        seeds) and apply Wilcoxon signed-rank tests with Holm correction.
  \item Sweep Late Acceptance list length over $\{1, 5, 50, 100, 500, 1000\}$
        and population size over $\{5, 10, 20, 50\}$.
\end{enumerate}

\section*{References}

\begin{description}
\setlength\itemsep{0pt}
\item Burke, E.~K. and Bykov, Y. (2017). The late acceptance hill-climbing
heuristic. \textit{European Journal of Operational Research}, 258(1), 70--78.

\item Burke, E.~K., Gendreau, M., Hyde, M., Kendall, G., Ochoa, G., \"{O}zcan,
E. and Qu, R. (2013). Hyper-heuristics: a survey of the state of the art.
\textit{Journal of the Operational Research Society}, 64(12), 1695--1724.

\item Cowling, P., Kendall, G. and Soubeiga, E. (2001). A hyperheuristic approach
to scheduling a sales summit. In \textit{Practice and Theory of Automated
Timetabling III}, LNCS 2079. Springer, 176--190.

\item Drake, J.~H., \"{O}zcan, E. and Burke, E.~K. (2012). An improved choice
function heuristic selection for cross domain heuristic search. In \textit{Parallel
Problem Solving from Nature}, PPSN XII, LNCS 7492. Springer, 307--316.

\item Friedman, J.~H. (2001). Greedy function approximation: a gradient boosting
machine. \textit{The Annals of Statistics}, 29(5), 1189--1232.

\item Lundberg, S.~M. and Lee, S.-I. (2017). A unified approach to interpreting
model predictions. In \textit{Advances in Neural Information Processing Systems
30}, 4765--4774.

\item Ochoa, G., Hyde, M., Curtois, T., Vazquez-Rodriguez, J.~A., Walker, J.,
Gendreau, M., Kendall, G., McCollum, B., Parkes, A.~J., Petrovic, S. and Burke,
E.~K. (2012). HyFlex: a benchmark framework for cross-domain heuristic search. In
\textit{Evolutionary Computation in Combinatorial Optimization}, LNCS 7245. Springer,
136--147.

\item Pillay, N. and Qu, R. (2018). \textit{Hyper-Heuristics: Theory and
Applications}. Springer. Prescribed textbook.

\item Ribeiro, M.~T., Singh, S. and Guestrin, C. (2016). ``Why should I trust
you?'': explaining the predictions of any classifier. In \textit{Proceedings of the
22nd ACM SIGKDD International Conference on Knowledge Discovery and Data Mining},
1135--1144.

\item Yates, W.~B., Keedwell, E.~C. and Kheiri, A. (2025). Explainable
optimisation through online and offline hyper-heuristics. \textit{ACM Transactions
on Evolutionary Learning and Optimization}, 5(2), article 10, 1--29.
\end{description}

\end{document}
